\section{在 APOGEE 上的应用 Application to APOGEE}\label{sec:apogee}

我们现在问：同一套计算能否直接拟合约减后的巡天光谱。我们把物理拟合与保留的标准星谱线定标结合，应用到 APOGEE 第 14 次数据释放（DR14）。这台光纤馈源近红外光谱仪及其交付分辨率见 \citet{Wilson2019}。选用该释放是因为它提供了受控比较所需的合并光谱、仪器定标与巨星分析产物 \citep{Abolfathi2018,Holtzman2018}。APOGEE 恒星参数与化学丰度流水线（ASPCAP）的标签提供外部对照 \citep{GarciaPerez2016,Jonsson2018}。

\begin{EN}
We now ask whether the same calculation can fit reduced survey spectra
directly. We combine physical fitting with the retained standard-star line
calibration and apply it to APOGEE Data Release 14 (DR14).
The fiber-fed near-infrared spectrograph and its delivered resolution are
described by \citet{Wilson2019}. We use this release because it provides the combined spectra, instrument
calibrations, and giant-star analysis products needed for the controlled
comparison \citep{Abolfathi2018,Holtzman2018}. Labels from the APOGEE Stellar
Parameters and Chemical Abundances Pipeline (ASPCAP) provide the
external comparison
\citep{GarciaPerez2016,Jonsson2018}.
\end{EN}

\begin{physbox}{APOGEE 与 ASPCAP 背景}
APOGEE 是架设在阿帕奇点天文台的高分辨率（$R\sim22{,}500$）近红外（1.5--1.7 $\mu$m，H 波段）巡天，观测了数十万颗银河系恒星，以红巨星为主。近红外的优势是尘埃消光小，能看穿银河盘。ASPCAP 是其官方分析流水线：把观测光谱与预计算的合成光谱网格（由 FERRE 驱动）做多维 $\chi^2$ 匹配，输出恒星参数与二十余种元素丰度，最后再施加经验零点与温度定标。本文的做法与 ASPCAP 的根本区别：不预计算网格、不插值，每颗星都从物理模型实时合成；也不施加 ASPCAP 的经验零点——这正是后文 Si 丰度出现刚性偏移的原因。
\end{physbox}

在这一巡天演示中，DR14 把其原始光谱型 $T_{\rm eff}$ 定标到低消光光度测量标度，把巨星 $\log g$ 定标到星震学重力，并把微观湍流作为单独的巨星坐标拟合 \citep{Holtzman2018}。我们不推导新的光谱—光度温度关系；本概念验证保留已发表的 $T_{\rm eff}$、$\log g$、$\xi$，以单独考察多元素与仪器问题。生产级分析可以把光谱与宽波段能谱分布、视差联合拟合 \citep{Cargile2020}。附录~\ref{app:stellar_label_sensitivity} 重复了放开 $T_{\rm eff}$、$\log g$、$\xi$ 的合成搜索，检验巨星支、红团簇与丰度形态对该选择的敏感性。

\begin{EN}
For the survey demonstration, DR14 calibrates its raw spectroscopic $T_{\rm eff}$ to a
low-reddening photometric scale and giant $\log g$ to asteroseismic gravities
and fits microturbulence as a separate giant-star coordinate
\citep{Holtzman2018}. We do not derive a new relation between spectroscopic and
photometric temperature. For this proof of concept, we retain the published
$T_{\rm eff}$, $\log g$, and $\xi$ to isolate the many-element and instrument
problem. A production analysis could instead fit the spectrum jointly with a
broadband spectral energy distribution and parallax \citep{Cargile2020}.
Appendix~\ref{app:stellar_label_sensitivity} repeats the synthesis
search with $T_{\rm eff}$, $\log g$, and $\xi$ free to test the sensitivity of
the giant-branch, red-clump, and abundance morphology to this choice.
\end{EN}

我们拟合 $[\mathrm{Fe}/\mathrm{H}]$，C、N、O、Mg、Al、Si、S、Ca、Ti、Mn、Ni 的 $[\mathrm{X}/\mathrm{Fe}]$，残余速度 $v_r$ 与有效高斯致宽 $v_b$。每个保留元素满足 $[\mathrm{X}/\mathrm{H}]=[\mathrm{Fe}/\mathrm{H}]+[\mathrm{X}/\mathrm{Fe}]$，未选金属跟随 $[\mathrm{Fe}/\mathrm{H}]$。$v_b$ 合并自转、宏观致宽以及平均仪器模型未表示的残余轮廓宽度。

\begin{EN}
We fit $[\mathrm{Fe}/\mathrm{H}]$, $[\mathrm{X}/\mathrm{Fe}]$ for C, N, O,
Mg, Al, Si, S, Ca, Ti, Mn, and Ni, residual velocity $v_r$, and effective
Gaussian broadening $v_b$. For every retained element,
$[\mathrm{X}/\mathrm{H}]=[\mathrm{Fe}/\mathrm{H}]+[\mathrm{X}/\mathrm{Fe}]$,
while unselected metals follow $[\mathrm{Fe}/\mathrm{H}]$. The fitted width
$v_b$ combines rotation, macroscopic broadening, and residual profile width not
represented by the mean instrument model.
\end{EN}

合成返回原生网格上的总量与连续谱流量。归一化之前，先把两者映射到观测像素：
\begin{equation}
f_p(\boldsymbol{\psi})=
\frac{\left\{\mathcal{K}\mathcal{S}(v_r)\mathcal{B}(v_b)
F_\lambda(\boldsymbol{\theta})\right\}_p}
{\left\{\mathcal{K}\mathcal{S}(v_r)\mathcal{B}(v_b)
F_{\lambda,{\rm c}}(\boldsymbol{\theta})\right\}_p},
\qquad
\boldsymbol{\psi}=(\boldsymbol{\theta},v_r,v_b).
\end{equation}
\noindent 其中 $\mathcal{B}$ 作用于对数合成网格，$\mathcal{S}$ 施加多普勒平移，$\mathcal{K}$ 施加随波长变化的线扩展函数与 APOGEE 原生采样。先对两者卷积再相除，保留了非线性的归一化。

\begin{EN}
Synthesis returns native-grid total and continuum fluxes. We map both to the
observed pixels before normalization,
\begin{equation}
f_p(\boldsymbol{\psi})=
\frac{\left\{\mathcal{K}\mathcal{S}(v_r)\mathcal{B}(v_b)
F_\lambda(\boldsymbol{\theta})\right\}_p}
{\left\{\mathcal{K}\mathcal{S}(v_r)\mathcal{B}(v_b)
F_{\lambda,{\rm c}}(\boldsymbol{\theta})\right\}_p},
\qquad
\boldsymbol{\psi}=(\boldsymbol{\theta},v_r,v_b).
\label{eq:apogee_forward}
\end{equation}
\noindent where $\mathcal{B}$ acts on the logarithmic synthesis grid,
$\mathcal{S}$ applies the Doppler shift, and $\mathcal{K}$ applies the
wavelength-dependent line-spread function and native APOGEE sampling.
Convolving both before division preserves the nonlinear normalization.
\end{EN}

本实验使用 DR14 定标导出的 LSF \citep{Nidever2015,Holtzman2018}，对六根代表性光纤取平均。一次性构建把从 $R_{\rm grid}=300{,}000$ 合成网格的插值折叠成紧凑的波长依赖权重并常驻 GPU；该算子直接映射到 APOGEE 像素，无需构建稠密卷积矩阵。换用光纤、曝光或探测器专属的 LSF 只需更换一份预置 GPU 资产。致宽、平移、LSF 卷积与采样在 H100 上对每对“总量—连续”流量约 $240~\mu\mathrm{s}$，而原生合成为 $1.4$ 秒——仪器模型的额外成本可忽略。公开拟合器还提供 GPU 常驻自转 kernel，本演示保留单一高斯宽度。

\begin{EN}
For this experiment, we use a DR14 calibration-derived LSF
\citep{Nidever2015,Holtzman2018}, averaged over six representative fibers. A
one-time build folds interpolation from the $R_{\rm grid}=300{,}000$ synthesis
grid into compact wavelength-dependent weights that remain on the GPU. The
operator maps directly onto the APOGEE pixels without constructing a dense
convolution matrix. A fiber-, exposure-, or detector-specific LSF requires
only a different prepared GPU asset. Broadening, shifting, LSF convolution,
and sampling take about $240~\mu\mathrm{s}$ per total--continuum flux pair on
an H100, compared with $1.4$~s for native synthesis. The instrument model
therefore adds negligible cost. The public fitter also provides a GPU-resident
rotational kernel, although this demonstration retains one Gaussian width.
\end{EN}

伪连续谱归一化的 DR14 光谱中仍残留平滑的连续谱差异 \citep{Holtzman2018}。每次非线性试探时，加权线性最小二乘在每个探测器上独立轮廓化一组四阶勒让德基 \citep{GolubPereyra2003}。似然使用已发表的不确定度与掩模 \citep{Holtzman2018}，外加 0.5\% 的模型下限。由于合并光谱组装在公共静止波长网格上 \citep{Nidever2015}，拟合从 $v_r=0$ 出发，只解局部残余平移。

\begin{EN}
Smooth continuum differences remain in the pseudo-continuum-normalized DR14
spectra \citep{Holtzman2018}. At every nonlinear trial, weighted linear least
squares profiles a fourth-order Legendre basis independently on each detector
\citep{GolubPereyra2003}. The likelihood uses the published uncertainties and
masks \citep{Holtzman2018} with a 0.5 percent model floor.
Because the combined spectra are assembled on a common rest-wavelength grid
\citep{Nidever2015}, the fit begins at $v_r=0$ and solves only a local
residual shift.
\end{EN}

观测光谱搜索由初始化大气与直接合成驱动。每增加一个丰度只需一条响应光谱，成本随标签数近似线性增长而非组合爆炸。随后物理求解器在所选混合物处收敛大气，一条新的 GPU 光谱提供任何最终的局部修正。被接受的模型永远使用物理收敛大气加直接合成。这一 12 丰度拟合在 H100 上搜索约 40 秒/星，大气收敛在多核 CPU 上再加约 30 秒。

\begin{EN}
Initialized atmospheres and direct synthesis drive the observed-spectrum
search.
Each additional abundance requires one response spectrum, so this cost grows
approximately linearly rather than combinatorially with the label count. The
physical solver then converges an atmosphere at the selected mixture, and a
fresh GPU spectrum supplies any final local correction. The accepted model
always uses a physically converged atmosphere and direct synthesis. For this
12-abundance fit, the H100 search takes about 40~s per star, and atmosphere
convergence adds about 30~s on multicore CPUs.
\end{EN}

作为概念验证，我们选取 1,600 颗巨星，在高、低 $\alpha$ APOGEE 序列与金属丰度上均衡采样 \citep{Nidever2014,Hayden2015}，并对族群轮廓重加权以代表母样本。图~\ref{fig:apogee_spectrum} 首先在像素空间检验计算：跨全部三个探测器，前向模型跟随谱线森林，轮廓化连续谱只吸收平滑的残余结构。这确立了仪器算子与多元素光谱可以联合拟合。

\begin{EN}
As a proof of concept, we select 1,600 giants, balanced across the high- and
low-$\alpha$ APOGEE sequences and metallicity
\citep{Nidever2014,Hayden2015}, and reweight the population contours to
represent the parent sample. Figure~\ref{fig:apogee_spectrum} first tests the
calculation in pixel space.
Across all three detectors, the forward model follows the line forest while the
profiled continuum absorbs only smooth residual structure. This establishes
that the instrument operators and many-element spectrum can be fitted together.
\end{EN}

随后检验独立拟合的丰度是否保持族群结构。图~\ref{fig:apogee_structure} 在被采样的金属丰度范围内保留了两条主要的 $[\mathrm{Mg}/\mathrm{Fe}]$ 序列，而不是塌缩成单一整体金属丰度趋势。图~\ref{fig:apogee_elements} 把比较推广到逐个元素：O、Si、S、Ca、Ti 在两个输入族群之间保持了化学分离，但绝对零点不同。

\begin{EN}
We then test whether the independently fitted abundances preserve population
structure. Figure~\ref{fig:apogee_structure} retains the two principal
$[\mathrm{Mg}/\mathrm{Fe}]$ sequences over the sampled metallicity range rather
than collapsing them into one bulk-metallicity trend. Figure~\ref{fig:apogee_elements}
extends the comparison element by element. O, Si, S, Ca, and Ti retain chemical
separation between the two input populations, although their absolute zero
points differ.
\end{EN}

\begin{physbox}{高/低 $\alpha$ 双序列为什么重要}
银河系的恒星在 $[\alpha/\mathrm{Fe}]$--$[\mathrm{Fe}/\mathrm{H}]$ 平面上分成两条序列：高 $\alpha$ 星形成早、增丰快（以 II 型超新星为主，$\alpha$ 元素多）；低 $\alpha$ 星形成晚、增丰慢（Ia 型超新星贡献了大量铁，压低了 $\alpha/$铁比）。这两条序列是银河系化学演化与并合历史的化石记录。一个丰度分析流水线最重要的检验之一就是：独立拟合能否重现这条双序列，而不是把它涂抹成单一趋势。本文的直接物理拟合做到了——说明 12 个丰度坐标各自携带了真实的天体物理信息，而非流水线伪影。Si 的刚性偏移则反映零点约定差异（DR14 做了经验零点修正而本文没有），属于标度问题而非形态错误。
\end{physbox}

近乎刚性的 Si 偏移与标度差异一致：DR14 对单元素丰度施加了经验温度与零点定标，Si 的调整幅度位居前列，而我们不做任何 APOGEE 零点修正 \citep{Holtzman2018}。Mn 向高金属丰度上升，与依赖金属丰度的产率及可观的 Ia 型超新星贡献一致；Ni 大体跟随 Fe \citep{Weinberg2019}。C 与 N 保持分别受约束，保留了 APOGEE 红巨星中随质量与演化变化的 C/N 信息 \citep{CharbonnelLagarde2010,Masseron2015,Martig2016}。最终的丰度修正通常在百分之几 dex 量级，不改变这一族群形态。

\begin{EN}
The nearly rigid Si offset is consistent with a scale difference because DR14 applies empirical
temperature and zero-point calibrations to individual abundances, with Si
among its larger adjustments, while we apply no APOGEE zero-point correction
\citep{Holtzman2018}. Mn rises toward higher metallicity, consistent with
metallicity-dependent yields and a
substantial Type Ia supernova contribution, while Ni broadly tracks Fe
\citep{Weinberg2019}. C and N remain separately constrained, preserving the
mass- and evolution-dependent C/N information measured in APOGEE red giants
\citep{CharbonnelLagarde2010,Masseron2015,Martig2016}. The final
abundance corrections are typically at the hundredth-dex level and leave this
population morphology unchanged.
\end{EN}

谱线表的迁移效果是独立可测的。在匹配的对照样本中，定标降低了每颗星的每有效像素平均卡方 $Q$：中位数从 3.32 降到 1.56，成对下降的中位数为 53\%。

\begin{EN}
The line-list transfer is independently measurable. In a matched control
sample, the calibration lowers the mean chi-square per valid pixel, $Q$, for
every star. The median $Q$ falls from
3.32 to 1.56, with a median paired reduction of 53 percent.
\end{EN}

定标谱线表在不强迫与 ASPCAP 丰度标度一致的情况下改善了像素级目标。剩余零点差异可来自大气物理、太阳混合物、谱线选择、掩模与连续谱处理。光谱拟合与逐元素序列共同表明：直接物理前向模型无需光谱模拟器即可恢复自洽的巡天化学。要建立定标后的丰度标度，则需要超出本概念验证的更大标准星集合与巡天专属的谱线整编。

\begin{EN}
The calibrated line list improves the pixel-level objective without forcing
agreement with the ASPCAP abundance scale. The remaining zero-point differences
can reflect the atmosphere physics, solar mixture, line selection, masks, and
continuum treatment. Together, the spectral fit and the element-by-element
sequences show that a direct physical forward model can recover coherent survey
chemistry without a spectral emulator. Establishing a calibrated abundance
scale would require a broader standard-star set and survey-specific line
curation beyond this proof of concept.
\end{EN}

\begin{figure}[t]
\centering
\includegraphics[width=0.98\textwidth]{figure_11.pdf}
\caption{代表性的 APOGEE DR14 拟合（收敛大气计算之后）。黑色为观测归一化光谱，橙虚线为经速度平移、致宽、线扩展函数卷积与像素采样后的轮廓化 \paynezero\ 模型，灰带为采用的 $1\sigma$ 不确定度。上光谱与残差面板覆盖全部三个探测器；下三对放大每个探测器的一段富线区间。残差以像素不确定度为单位。（Representative APOGEE DR14 fit.）}
\label{fig:apogee_spectrum}
\end{figure}

\begin{figure}[t]
\centering
\includegraphics[width=0.98\textwidth]{figure_12.pdf}
\caption{均衡 1,600 目标 APOGEE 样本的化学结构。前两图为 DR14 定标星表与 \paynezero\ 拟合的 $[\mathrm{Mg}/\mathrm{Fe}]$--$[\mathrm{Fe}/\mathrm{H}]$，橙、蓝标记 APOGEE 高、低 $\alpha$ 样本；第三图比较两者的铁丰度，虚线为相等线。等值线包围逆选择权重密度的 50、80、95\%。直接拟合保留了两条化学序列，同时允许丰度零点平移。（Chemical structure of the balanced 1,600-target APOGEE sample.）}
\label{fig:apogee_structure}
\end{figure}

\begin{figure}[t]
\centering
\includegraphics[width=0.98\textwidth]{figure_13.pdf}
\caption{均衡 1,600 目标 \paynezero\ 拟合的逐元素丰度分布。每面板为一个 $[\mathrm{X}/\mathrm{Fe}]$ 对 $[\mathrm{Fe}/\mathrm{H}]$；橙、蓝点与等值线分别来自高、低 $\alpha$ 样本；灰虚等值线为同一批恒星与权重下 APOGEE DR14 的 50、80\% 密度。拟合的丰度平面呈现连贯的、随元素变化的序列，而非单一整体金属丰度坐标。（Elemental-abundance distributions from the balanced 1,600-target fit.）}
\label{fig:apogee_elements}
\end{figure}
